\textbf{Saturn and Titan.} The official answer depends on the night of
observation; reference values are given for both nights. Titan's position is
checked with a transparent template foil: if Titan is within the template
circle, 2 points; if outside, 0 points.

% FIGURE NEEDED: reference view-field drawings with Saturn, the labelled
% moons and the cross marking Titan (T), plus the reference ring drawing, for
% 2019.08.04 22:00 UT and 2019.08.05 22:00 UT, and the transparent template
% foils (2019_TEL_sol.pdf pages 1-4)

\medskip
For 2019.08.04 22:00 UT:
\[ d_{\mathrm{Titan}} = 164'' \]
\begin{center}
\begin{tabular}{ll}
$d$: $148$--$180''$ & 3 points \\
$d$: $139$--$189''$ & 2 points \\
$d$: $123$--$205''$ & 1 point \\
\end{tabular}
\end{center}
\[ \mathrm{PA}_{\mathrm{Titan}} = 109^\circ \]
\begin{center}
\begin{tabular}{ll}
PA: $99$--$119^\circ$ & 4 points \\
PA: $89$--$129^\circ$ & 2 points \\
\end{tabular}
\end{center}

For 2019.08.05 22:00 UT:
\[ d_{\mathrm{Titan}} = 129'' \]
\begin{center}
\begin{tabular}{ll}
$d$: $116$--$131''$ & 3 points \\
$d$: $110$--$148''$ & 2 points \\
$d$: $97$--$161''$ & 1 point \\
\end{tabular}
\end{center}
\[ \mathrm{PA}_{\mathrm{Titan}} = 124^\circ \]
\begin{center}
\begin{tabular}{ll}
PA: $114$--$134^\circ$ & 4 points \\
PA: $104$--$144^\circ$ & 2 points \\
\end{tabular}
\end{center}

\medskip
Scoring:

\begin{center}
\begin{tabular}{lll}
Titan position + name: & & 2 points \\
Titan is outside the circle: & & 0 point \\[4pt]
Ring continuity in front of the disk is proper & & 2 points \\
Ring tilt towards Earth is proper & & 2 points \\
\multicolumn{3}{l}{\quad ring does not exceed northern pole, ring is not too thin} \\
Ring tilt towards sky E--W direction is proper & & 2 points \\[4pt]
PA measurement of Titan: & $\pm 10^\circ$ & 4 points \\
 & $\pm 20^\circ$ & 2 points \\[4pt]
Distance calculation of Titan & $\pm 10$\% & 3 points \\
 & $\pm 15$\% & 2 points \\
 & $\pm 25$\% & 1 point \\
\end{tabular}
\end{center}

The Saturn drawing evaluation scheme uses the sub-Earth latitude
$\Phi = 24.8^\circ$ (ring tilt towards Earth, max 2 points: drawings
equivalent to $22.8$--$26.8^\circ$ get 2p; $20.8^\circ$ gets 1p, S pole
visible; $28.8^\circ$ gets 1p, very thick; thinner or thicker gets 0p) and
the equatorial position angle $\mathrm{PA} = 6.4^\circ$ (ring tilt towards
the sky E--W direction, max 2 points: $6.4^\circ$ or $9.4^\circ$ get 2p;
$3.4^\circ$ or $12.4^\circ$ get 1p; $0^\circ$ or $15.4^\circ$ get 0p).
Inverse, improper, too wide or too narrow rings get 0 points for ring
continuity.

% FIGURE NEEDED: "Saturn drawing evaluation scheme" -- galleries of example
% ring drawings with their scores (2019_TEL_sol.pdf pages 5-6)

\medskip
\textbf{Alternative task: observing Epsilon Lyr.} \hfill TOTAL: 15 points

Checking the object in the FOV, and pointing is evaluated by the telescope
assistant. \hfill(0--3 points)

FOV with 10 mm eyepiece: the components of the close pairs are not resolved.

% FIGURE NEEDED: reference FOV drawing with epsilon-1 and epsilon-2 Lyr, the
% surrounding star field and the N/E direction arrows
% (2019_TEL_sol.pdf page 7)

FOV drawing: star field and directions:
\begin{center}
\begin{tabular}{ll}
Correct direction \& labelling of North and East relative to the star field: & 2 points \\
Correct drawing of at least 3 stars in the FOV: & 1 point \\
Correct drawing of at least 6 stars in the FOV: & 2 points \\
\end{tabular}
\end{center}

Wide pair distance: $208'' = 3.47'$ \hfill 2 points

\begin{center}
\begin{tabular}{lll}
Distance estimation of wide pair: & $d_{\epsilon 1 - \epsilon 2} = 3.2$--$3.8'$ & 2 points \\
 & $d_{\epsilon 1 - \epsilon 2} = 2.8$--$4.2'$ & 1 point \\
\end{tabular}
\end{center}

Wide pair PA: $172^\circ$ \hfill 3 points

\begin{center}
\begin{tabular}{lll}
Position angle of the wide pair: & $\mathrm{PA}_{\epsilon 2} = 170$--$175^\circ$ & 3 points \\
 & $\mathrm{PA}_{\epsilon 2} = 167$--$177^\circ$ & 2 points \\
 & $\mathrm{PA}_{\epsilon 2} = 162$--$182^\circ$ & 1 point \\
\end{tabular}
\end{center}

Correct estimation of the relative angle of close pairs: \hfill 3 points

Relative angle between the direction of lines fitted onto the two close
pairs: $92^\circ$ (also the $88^\circ$ complementary angle is acceptable ---
referring to this, the evaluation bands are centered to $90^\circ$):
\begin{center}
\begin{tabular}{ll}
$85$--$95^\circ$ & 3 points \\
$80$--$100^\circ$ & 2 points \\
$75$--$105^\circ$ & 1 point \\
\end{tabular}
\end{center}
